\documentclass[a4paper]{article}
\usepackage[pdftex]{graphicx}
\usepackage[utf8]{inputenc}
\usepackage{enumerate}
\usepackage{icomma}
\usepackage{amssymb}
\usepackage{amsmath}
\usepackage{tikz}
\usepackage{href-ul}
\hypersetup{
	colorlinks=true,
	linkcolor=blue,
	urlcolor=blue}
\usepackage{geometry}
\geometry{a4paper, top=15mm, left=15mm, right=15mm, bottom=15mm,
	headsep=10mm, footskip=12mm}

\begin{document}
	\begin{tikzpicture}(10,3)
		\draw[ultra thick](2,0) --(10,0) -- (10,1.3) --(2,1.3) -- (2,0);
		\draw[fill=black](2,0)-- (0,.6) -- (2,.6) -- (2,0);
		\node at (6,.6) {\large Solution sheet from \href{http://www.okuyakl.com}{www.okuyakl.com}};
	\end{tikzpicture}
	\vspace{0.5 cm}
	
	\noindent{\bf Task 1.}\\
	First we solve the following inequality for the expression between the absolute lines:
	
	$$
	\begin{array}{rl}
		{4 \over 3}x-4 & \ge 0\\
		{4 \over 3}x & \ge 4\\
		x & \ge 3
	\end{array}
	$$
	\noindent case distinction:
	
	\underbar{Case 1: $x \ge 3$}\\
	$$f(x)=-{2 \over 3} x + {4 \over 3}x-4={2 \over 3}x -4 $$
	
	\underbar{Case 2: $x < 3$}\\
	$$f(x)=-{2 \over 3} x -{4 \over 3} x + 4 = -2 x +4 $$
	
	\noindent Now we get the following function defined section by section:
	
	\begin{minipage}{0.4\textwidth}
		$$
		f(x)=
		\begin{cases}
			-2x+4 & \text{for} \quad x < 3 \\
			{2 \over 3}x-4 & \text{for} \quad x \ge 3
		\end{cases}
		$$
	\end{minipage}
	\hfill
	\begin{minipage}{0.55\textwidth}
		\setlength{\unitlength}{0.5 cm}
		\begin{picture}(10,10)
			\put(0,3){\vector(1,0){10}}
			\put(1,0){\vector(0,1){10}}
			\put(4,1){\line(-1,2){4}}
			\put(4,1){\line(3,2){5}}
			\put(3.9,2.4){3}
			\put(0.4,6.8){4}
			\put(0.4,2.4){0}
			\put(0.4,9){y}
			\put(9,2.4){x}
			\put(9,4){$f(x)$}
		\end{picture}
	\end{minipage}
	
	\noindent{\bf Task 2.}\\
	$$ f_1 \Rightarrow E \quad f_2 \Rightarrow C \quad f_3 \Rightarrow B \quad f_4 \Rightarrow A \quad f_5 \Rightarrow F \quad f_6 \Rightarrow D $$
	\vspace{0.5 cm}
	
	\noindent
	\begin{minipage}{0.5\textwidth}
		\noindent{\bf Task 3. a) c)}\\
		\includegraphics[width=7 cm]{xvier032}\\
		\setlength{\unitlength}{0.5 cm}
		\begin{picture}(0,0)
			\put(4,9.5){$f^{-1}(x)$}
			\put(9,4){$f(x)$}
		\end{picture}
		
	\end{minipage}
	\hfill
	\begin{minipage}{0.5\textwidth}
		\noindent{\bf Task 3. b)}\\
		It is reversible because each y-value only occurs once, i.e. it is monotonic.
		
		\noindent{\bf Task 3. d)}\\
		We resolve:
		$$
		\renewcommand{\arraystretch}{2}
		\begin{array}{rcll}
			y &=& (x-3)^4+1 &|-1\\
			y-1 &=& (x-3)^4 &|\sqrt[4]{\quad}\\
			\sqrt[4]{y-1} &=& x-3 &|+3 \\
			\sqrt[4]{y-1} +3 &=& x &| x \leftrightarrow y\\
			f^{-1}(x)&=& \sqrt[4]{x-1} +3
		\end{array}
		$$
	\end{minipage}
	$$ \mathbb{D}_{f^{-1}}=]1;\infty[ \qquad \mathbb{W}_{f^{-1}}=]3;\infty[$$
	
	\noindent{\bf Task 4. a)}\\
	Asymptote $y=3 \Rightarrow$ $e=3$\\
	$P \in h_2 \Rightarrow$
	$$
	\renewcommand{\arraystretch}{2}
	\begin{array}{rcl}
		a(1-1)^{-2}+3 &=& 3.5 \\
		a +3 &=& 3.5 \\
		a &=& 0.5\\
		\Rightarrow &=& y=0.5(x-1)^{-1}+3
	\end{array}
	$$
	\noindent{\bf Task 4. b)}\\
	To show: $f(-x)=f(x)$
	$$(-x)^{-2}=\left((-x)^{2}\right)^{-1}=(x^{2})^{-1}= x^{-2 } \quad (w)$$
	The function is axisymmetric.
	
	\noindent{\bf Task 5. a)}\\
	We equate:
	$$
	\renewcommand{\arraystretch}{2}
	\begin{array}{rcll}
		(x-1)^{-1} + 3 &=& -4x +3 \\
		(x-1)^{-1} &=& -4x &|\cdot (x-1) \\
		1 &=& -4x^2 + 4x \\
		-4x^2 + 4x -1 &=& 0 \\
		x_{1/2}&=& {-4 \pm \sqrt{16 - 4 \cdot 4}\over 2\cdot (-4)}\\
		x &=& {1\over 2}
	\end{array}
	$$
	Inserting into the straight line equation:
	$$y= -4\cdot {1\over 2} + 3 = 1 \quad \Rightarrow \quad A(0.5|1)$$
	
	\noindent {\bf Task 5. b)}\\
	The x coordinate of point B satisfies the following equation:
	$$
	\renewcommand{\arraystretch}{1.5}
	\begin{array}{rcl}
		{x_B + x_A \over 2} &=& x_M \\
		x_B &=& 2 x_M - x_A \\
		x_B &=& 2 \cdot 1.75 - 0.5 \\
		x_B &=& 3 \\
		B \in g &\Rightarrow& y_B = -4 \cdot 3 +3\\
		y_B &=& -9 \quad \Rightarrow \quad B(3|-9)
	\end{array}
	$$
	
	\noindent {\bf Task 6. }\\
	\begin{tabular}{ccccccc}
		$ f(x)$&=& axisymmetric & \parbox[20 pt][2em][c]{0 cm}{} &
		$g(x)$&=& point symmetric\\
		$ h(x)$&=&axisymmetric & \parbox[20 pt][2em][c]{0 cm}{} &
		$i(x)$&=&point symmetric\\
		$ j(x)$&=& not symmetrical & \parbox[20 pt][2em][c]{0 cm}{} &
		$k(x)$&=& not symmetrical\\
	\end{tabular}
	
	\noindent {\bf Task 7.}\\
	\begin{enumerate}[(i)]
		\item n odd, a positive
		\item n even, a = 1
	\end{enumerate}
	
	\begin{center}
		\includegraphics[width=7 cm]{../../viecher/eendcomic.pdf}
		
		Here you can go back to the \href{https://www.okuyakl.de/math/m10fuaL032/ae032.pdf}{task sheet}
	\end{center}
	
\end{document}